\section{What the Table Shows, and Why}

Table~\ref{tab:results} suggests several rules. In this section we state them
one by one and prove each of them, so that they hold for every number of
people and not only for the sixteen circles we played. We first prove a lemma
that we use in all the proofs: the game only depends on the places of the
people, not on their numbers. The first two rules then link a circle of $2k$
or $2k+1$ people to a circle of $k$ people. The third one is about the way
each $n$ is written in the table.

\subsection{The key lemma: only the places matter}\label{sec:tool}

At any moment of the game, the people still present stand around the circle
in a fixed order. We write them as a list, starting with the person who is
skipped next.

\begin{defi}\label{def:list}
A \emph{list} $(a_1,a_2,\dots,a_k)$ is the sequence of the $k$ people still
present, in the order of the circle, where $a_1$ is the next person to be
skipped. If $k\geq2$, one \emph{step} of the game skips $a_1$ and eliminates
$a_2$. The next person to be skipped is then $a_3$, and $a_1$ is now the last
person before we come back to $a_3$:
\[
(a_1,a_2,a_3,\dots,a_k)\longrightarrow(a_3,a_4,\dots,a_k,a_1).
\]
For $k=2$, the new list is $(a_1)$. When $k=1$ the game is over and $a_1$ is
the survivor.
\end{defi}

Each step removes one person, so a game with $k$ people ends after $k-1$
steps. The game of Definition~\ref{def:game} starts from the list
$(1,2,\dots,n)$. For example, with $n=5$:
\[
(1,2,3,4,5)\to(3,4,5,1)\to(5,1,3)\to(3,5)\to(3).
\]
The people eliminated are $2,4,1,5$, in this order, and the survivor is $3$,
as in Table~\ref{tab:results}.

The step rule only uses the places in the list, never the numbers of the
people. This is what the next lemma says.

\begin{lemma}[Only the places matter]\label{lem:places}
Let $(a_1,\dots,a_k)$ be a list of $k\geq1$ people. The survivor of the game
that starts from this list is $a_{J(k)}$, the person at place $J(k)$.
\end{lemma}

\begin{proof}
We prove by induction on $k\geq1$ the statement
\[
L(k):\quad\text{for every list $(a_1,\dots,a_k)$ of $k$ people, the survivor
is $a_{J(k)}$.}
\]
\emph{Base case.} For $k=1$, the only person $a_1$ is the survivor, and
$J(1)=1$. So $L(1)$ is true.

\emph{Induction step.} Let $k\geq2$ and assume $L(k-1)$. Take a list
$(a_1,\dots,a_k)$. One step gives the list of $k-1$ people
\[
(b_1,\dots,b_{k-1})=(a_3,a_4,\dots,a_k,a_1),
\]
that is, $b_j=a_{f(j)}$ with
\[
f(j)=\begin{cases} j+2 & \text{if } 1\leq j\leq k-2,\\ 1 & \text{if } j=k-1.
\end{cases}
\]
The game goes on from this new list, so by $L(k-1)$ its survivor is
$b_{J(k-1)}=a_{f(J(k-1))}$. This is true for every list. For the list
$(1,2,\dots,k)$, which is the usual game with $k$ people, the survivor is
$J(k)$ by definition, so $J(k)=f\big(J(k-1)\big)$. Going back to any list, the
survivor is therefore $a_{f(J(k-1))}=a_{J(k)}$, and $L(k)$ is true.

By induction, $L(k)$ is true for every $k\geq1$.
\end{proof}

The proof also gives a relation between $J(k)$ and $J(k-1)$. Since
$J(k)=f\big(J(k-1)\big)$,
\begin{equation}\label{eq:onestep}
J(k)=\begin{cases} J(k-1)+2 & \text{if } J(k-1)\leq k-2,\\
1 & \text{if } J(k-1)=k-1 .\end{cases}
\end{equation}
In words: adding one person moves the survivor two places, except when the
survivor was the last person, in which case it goes back to person $1$. This
is what we saw in Table~\ref{tab:results}. Relation~\eqref{eq:onestep} gives
$J(n)$ for every $n$, but it needs $n-1$ steps. The two relations of the next
subsections are much faster, because they divide $n$ by $2$.

\subsection{Even circles: \texorpdfstring{$J(2k)=2J(k)-1$}{J(2k) = 2J(k) - 1}}\label{sec:even}

\subsubsection{What we observe}

We put side by side the survivor for $k$ people and the survivor for $2k$
people.
\begin{center}
\renewcommand{\arraystretch}{1.2}
\begin{tabular}{|c|c|c|c|c|c|c|c|c|}
\hline
$k$ & 1 & 2 & 3 & 4 & 5 & 6 & 7 & 8\\ \hline
$J(k)$ & 1 & 1 & 3 & 1 & 3 & 5 & 7 & 1\\ \hline
$J(2k)$ & 1 & 1 & 5 & 1 & 5 & 9 & 13 & 1\\ \hline
$2J(k)-1$ & 1 & 1 & 5 & 1 & 5 & 9 & 13 & 1\\ \hline
\end{tabular}
\end{center}
The last two lines are equal in all eight cases.

\begin{obs}\label{obs:even}
For every $k\geq1$, \ $J(2k)=2J(k)-1$.
\end{obs}

\subsubsection{Proof}

Let $k\geq1$, and let $n=2k$ people stand around the circle. The idea is to
play the first round completely, and to see that what is left is a game with
$k$ people.

\textbf{Step 1: the first round eliminates exactly the even numbers.} We show
by induction on $j$, for $j=1,\dots,k$, the statement
\[
E(j):\quad\text{the first $j$ eliminations skip $1,3,\dots,2j-1$ and
eliminate $2,4,\dots,2j$.}
\]
\emph{Base case.} The game starts by skipping person $1$ and eliminating
person $2$, so $E(1)$ is true.

\emph{Induction step.} Let $1\leq j<k$ and assume $E(j)$. At this point the
people eliminated are $2,4,\dots,2j$, and the last one is $2j$. The next
person is $2j+1$. The people
$2j+1$ and $2j+2$ are still present, and $2j+2\leq2k$. So the game skips
$2j+1$ and eliminates $2j+2$, and $E(j+1)$ is true.

By $E(k)$, the $k$-th person eliminated is $2k$, the last person of the circle. The
next person is person $1$, who is still present: we are back at the beginning
of the circle and the first round is over. It has eliminated
$2,4,\dots,2k$, and the $k$ people left are, in the order of the circle,
\[
1,\ 3,\ 5,\ \dots,\ 2k-1 .
\]
The game goes on by skipping person $1$.

\textbf{Step 2: the people left form a game with $k$ people.} We write the
people left as a list $(a_1,a_2,\dots,a_k)$, in the order of the circle,
starting with the person who is skipped next:
\[
a_i=2i-1\qquad(i=1,\dots,k),
\]
so $a_1=1$, $a_2=3$, \dots, $a_k=2k-1$ (Figure~\ref{fig:even}). The rest of the
game is the game of Definition~\ref{def:game} played on $a_1,\dots,a_k$ instead
of $1,\dots,k$: the circle has $k$ people and it starts by skipping the first
one.

\textbf{Step 3: the survivor is $a_{J(k)}$.} By Lemma~\ref{lem:places},
applied to the list $(a_1,\dots,a_k)$, the survivor is the person at place
$J(k)$ of this list. So
\[
J(2k)=a_{J(k)}=2J(k)-1 .
\]
This proves Observation~\ref{obs:even} for every $k\geq1$. \hfill$\square$

\begin{figure}[!ht]
\centering
\renewcommand{\nodesize}{7mm}\renewcommand{\nodefont}{\footnotesize}\footnotesize
\begin{tikzpicture}
  \draw[gray,dotted,very thick] (-35:2.1) arc (-35:-235:2.1);
  \node[person] (a1) at (90:2.1) {1};
  \node[gone] (a2) at (55:2.1) {}; \cross{a2} \node[fill=white,inner sep=0.4pt] at (a2) {2};
  \node[person] (a3) at (20:2.1) {3};
  \node[gone] (a4) at (-15:2.1) {}; \cross{a4} \node[fill=white,inner sep=0.4pt] at (a4) {4};
  \node[gone] (a6) at (125:2.1) {}; \cross{a6} \node[font=\scriptsize,fill=white,inner sep=0.4pt] at (a6) {$2k$};
  \node[person,font=\scriptsize,minimum size=8.5mm] at (160:2.1) {$2k{-}1$};
  \node[font=\scriptsize] at (90:2.75) {$a_1$};
  \node[font=\scriptsize] at (20:2.75) {$a_2$};
  \node[font=\scriptsize] at (160:2.8) {$a_k$};
\end{tikzpicture}
\caption{The circle of $2k$ people after the first round: $a_i=2i-1$. The
dotted arc stands for the people not drawn.}
\label{fig:even}
\end{figure}

\textbf{Example.} For $n=10$, the first round leaves $(1,3,5,7,9)$. Since
$J(5)=3$, the survivor is the third person of this list, $a_3=5$. So
$J(10)=2\times3-1=5$, as in Table~\ref{tab:results}.

\subsection{Odd circles: \texorpdfstring{$J(2k+1)=2J(k)+1$}{J(2k+1) = 2J(k) + 1}}

\subsubsection{What we observe}

We now compare the survivor for $k$ people and for $2k+1$ people.
\begin{center}
\renewcommand{\arraystretch}{1.2}
\begin{tabular}{|c|c|c|c|c|c|c|c|}
\hline
$k$ & 1 & 2 & 3 & 4 & 5 & 6 & 7\\ \hline
$J(k)$ & 1 & 1 & 3 & 1 & 3 & 5 & 7\\ \hline
$J(2k+1)$ & 3 & 3 & 7 & 3 & 7 & 11 & 15\\ \hline
$2J(k)+1$ & 3 & 3 & 7 & 3 & 7 & 11 & 15\\ \hline
\end{tabular}
\end{center}
Again, the last two lines are equal in all seven cases.

\begin{obs}\label{obs:odd}
For every $k\geq1$, \ $J(2k+1)=2J(k)+1$.
\end{obs}

\subsubsection{Proof}

Let $k\geq1$, and let $n=2k+1$ people stand around the circle. We follow the
same three steps as for $2k$ people. Only the end of the first round changes.

\textbf{Step 1: the first round.} The induction of the previous proof works
without change: the first $k$ eliminations skip $1,3,\dots,2k-1$ and
eliminate $2,4,\dots,2k$. Indeed, for $j<k$ we still have $2j+2\leq2k<2k+1$.

After the $k$-th elimination (person $2k$), one person remains before the end
of the circle: person $2k+1$. He is skipped. The next person is person $1$,
since after person $2k+1$ we come back to the beginning of the circle, and
person $1$ is eliminated. As said in Section~\ref{sec:conventions}, this elimination belongs to
the first round. So the first round eliminates $2,4,\dots,2k$ and then $1$,
and the $k$ people left are
\[
3,\ 5,\ 7,\ \dots,\ 2k+1 .
\]
The game goes on with the first person after $1$ who is still present.
Person $2$ is already eliminated, so this is person $3$, and he is the next one
skipped.

\textbf{Step 2: the people left form a game with $k$ people.} This time the
list starts with person $3$:
\[
a_i=2i+1\qquad(i=1,\dots,k),
\]
so $a_1=3$, $a_2=5$, \dots, $a_k=2k+1$ (Figure~\ref{fig:odd}). The rest of the
game is again the game of Definition~\ref{def:game} with $k$ people, played on
$a_1,\dots,a_k$.

\textbf{Step 3: the survivor.} By Lemma~\ref{lem:places}, the survivor is
the person at place $J(k)$ of the list, so
\[
J(2k+1)=a_{J(k)}=2J(k)+1 .
\]
This proves Observation~\ref{obs:odd} for every $k\geq1$. \hfill$\square$

\begin{figure}[!ht]
\centering
\renewcommand{\nodesize}{7mm}\renewcommand{\nodefont}{\footnotesize}\footnotesize
\begin{tikzpicture}
  \draw[gray,dotted,very thick] (-35:2.1) arc (-35:-235:2.1);
  \node[gone] (b1) at (90:2.1) {}; \cross{b1} \node[fill=white,inner sep=0.4pt] at (b1) {1};
  \node[gone] (b2) at (55:2.1) {}; \cross{b2} \node[fill=white,inner sep=0.4pt] at (b2) {2};
  \node[person] (b3) at (20:2.1) {3};
  \node[gone] (b4) at (-15:2.1) {}; \cross{b4} \node[fill=white,inner sep=0.4pt] at (b4) {4};
  \node[gone] (b5) at (160:2.1) {}; \cross{b5} \node[font=\scriptsize,fill=white,inner sep=0.4pt] at (b5) {$2k$};
  \node[person,font=\scriptsize,minimum size=8.5mm] at (125:2.1) {$2k{+}1$};
  \node[font=\scriptsize] at (20:2.75) {$a_1$};
  \node[font=\scriptsize] at (125:2.8) {$a_k$};
\end{tikzpicture}
\caption{The circle of $2k+1$ people after the first round: $a_i=2i+1$.
Person $1$ is eliminated at the end of the round, so the list starts with
person $3$.}
\label{fig:odd}
\end{figure}

\textbf{Example.} For $n=11$, the first round eliminates $2,4,6,8,10$ and then
$1$, and leaves $(3,5,7,9,11)$. Since $J(5)=3$, the survivor is $a_3=7$. So
$J(11)=2\times3+1=7$, as in Table~\ref{tab:results}.

\FloatBarrier
\subsubsection{Computing \texorpdfstring{$J(n)$}{J(n)} by halving}\label{sub:halving}

With $J(1)=1$, Observations~\ref{obs:even} and~\ref{obs:odd} give $J(n)$ for
every $n$. We divide $n$ by $2$ (rounding down) until we reach $1$, then we go
back up with the two relations. For $n=100$:
\[
100\to50\to25\to12\to6\to3\to1,
\]
and
\begin{align*}
J(3)&=2J(1)+1=3, & J(6)&=2J(3)-1=5, & J(12)&=2J(6)-1=9,\\
J(25)&=2J(12)+1=19, & J(50)&=2J(25)-1=37, & J(100)&=2J(50)-1=73.
\end{align*}
For $n=1000$, the halvings are
$1000\to500\to250\to125\to62\to31\to15\to7\to3\to1$, and going back up gives
\begin{align*}
&J(3)=3, && J(7)=7, && J(15)=15, && J(31)=31, && J(62)=61,\\
&J(125)=123, && J(250)=245, && J(500)=489, && J(1000)=977 .
\end{align*}
About $\log_2 n$ steps replace the $n-1$ eliminations of the game: $6$ steps
for $n=100$ and $9$ steps for $n=1000$. Still, this does not tell us what $J(n)$
looks like. For that we go back to the way the numbers are written in
Table~\ref{tab:results}.

\subsection{Every \texorpdfstring{$n$}{n} is \texorpdfstring{$2^m+\ell$}{2\^{}m+l} in exactly one way}

\subsubsection{What we observe}

In Table~\ref{tab:results}, each $n$ is written as $2^m+\ell$, where $2^m$ is
the largest power of $2$ that is not bigger than $n$:
\[
5=2^2+1,\qquad 10=2^3+2,\qquad 15=2^3+7,\qquad 16=2^4+0 .
\]
In every line the remainder $\ell$ is smaller than $2^m$, and there is only
one such way to write $n$. For instance, $10=2^2+6$ is not allowed, because
$6\geq2^2$.

\begin{obs}\label{obs:decomp}
Every integer $n\geq1$ can be written in exactly one way as
\[
n=2^m+\ell,\qquad m\geq0,\quad 0\leq\ell<2^m .
\]
\end{obs}

\subsubsection{Proof by induction}

\textbf{Existence, by strong induction on $n$.} For $n=1$ we have
$1=2^0+0$, and $0<2^0$.

Let $n\geq2$, and assume that every integer $k$ with $1\leq k<n$ can be
written $k=2^p+r$ with $0\leq r<2^p$. We write $n=2k$ if $n$ is even, or
$n=2k+1$ if $n$ is odd. In both cases $1\leq k<n$, so $k=2^p+r$ with
$0\leq r<2^p$.
\begin{itemize}
  \item If $n=2k$, then $n=2^{p+1}+2r$, and $0\leq2r<2\cdot2^p=2^{p+1}$.
  \item If $n=2k+1$, then $n=2^{p+1}+(2r+1)$. Since $r\leq2^p-1$, we get
        $2r+1\leq2^{p+1}-1<2^{p+1}$.
\end{itemize}
In both cases $n=2^m+\ell$ with $m=p+1$ and $0\leq\ell<2^m$.

\textbf{Uniqueness.} The condition $0\leq\ell<2^m$, with $\ell=n-2^m$, means
\begin{equation}\label{eq:block}
2^m\leq n<2^{m+1},
\end{equation}
since $\ell\geq0$ means $2^m\leq n$, and $\ell<2^m$ means
$n<2^m+2^m=2^{m+1}$. Assume that $n=2^m+\ell=2^{m'}+\ell'$ are two such
writings, with $m<m'$. Then $m+1\leq m'$, and \eqref{eq:block} for $m$ and for
$m'$ gives
\[
n<2^{m+1}\leq2^{m'}\leq n,
\]
which is impossible. In the same way $m'<m$ is impossible. So $m=m'$, and
then $\ell=n-2^m=\ell'$. \hfill$\square$

The proof also shows how to find $m$ and $\ell$: by \eqref{eq:block}, $2^m$ is
the largest power of $2$ not bigger than $n$, and $\ell=n-2^m$. For example,
$64\leq100<128$ gives $100=2^6+36$. The integers with the same $m$ form the
blocks $2^m\leq n<2^{m+1}$ separated by lines in Table~\ref{tab:results}.

The three results of this section are the tools we need. In the next section
we use them to prove a formula for $J(n)$.
